\chapter{Anti-PLA2R Antibodies}

\section*{What Is This Test?}
Anti-phospholipase A2 receptor antibodies are autoantibodies associated with primary membranous nephropathy, a kidney disease that commonly presents with significant proteinuria or nephrotic syndrome.

\section*{Alternative Names}
Anti-PLA2R; PLA2R Antibody; Phospholipase A2 Receptor Antibody.

\section*{Why This Test Is Ordered}
Testing supports evaluation of adults with unexplained nephrotic-range proteinuria or suspected membranous nephropathy. It can help assess disease activity and response to treatment, but it does not replace kidney-function testing or specialist evaluation.

\section*{How This Test Is Performed}
A blood sample is analyzed by an immunoassay. Urine protein measurement, serum albumin, creatinine, eGFR, and evaluation for secondary causes are interpreted alongside the antibody result.

\section*{Preparation, Risks, and Limitations}
No fasting is required. Venipuncture is the usual risk. A negative result does not exclude membranous nephropathy, and a positive result does not rule out every secondary cause. Antibody levels may change with disease activity and therapy.

\section*{Understanding Results}
A positive result supports primary membranous nephropathy in the correct clinical context. Higher or persistent levels may be associated with active disease, while falling or negative levels can accompany immunologic remission; urine protein and kidney function remain essential for monitoring.

\section*{References}
Kidney Disease: Improving Global Outcomes guidance on glomerular diseases.

\clearpage
\section*{Understanding Results and Follow-up}
The laboratory report should identify the specimen, method, units, reference interval or decision limit, and whether the result is positive, negative, abnormal, or inconclusive. Reference intervals vary with age, sex, pregnancy, specimen type, assay, and laboratory; a result outside a range is not automatically a diagnosis, and a result within a range does not always exclude disease. Discuss unexpected results with the ordering clinician, who can decide whether repeat or confirmatory testing, treatment, or referral is appropriate. Do not change medicines or delay urgent care based on an isolated result.


\section*{Regulatory Status and Authorization Date}
This chapter describes a test category, not a specific commercial product. FDA, CDSCO, EU IVDR, and MHRA approval or registration dates therefore cannot be assigned without the assay manufacturer, product name, instrument, and intended use. For laboratory-developed tests, an individual agency approval date may not exist. See the Preface for the interpretation of regulatory status.

\clearpage
